\subsection{Herramientas y tecnologías}

A continuación, se describen las principales herramientas y tecnologías empleadas en el desarrollo del presente trabajo de investigación. Se detallan los lenguajes de programación, las librerías y \textit{frameworks} utilizados, así como los recursos computacionales necesarios.

En primer lugar, se destaca el uso de LaTeX para la redacción del presente documento, dado que permite estructurar adecuadamente la memoria, gestionar referencias bibliográficas y representar expresiones matemáticas de forma precisa.

Para la adquisición de datos se ha empleado la herramienta \textit{yfinance}, que permite acceder a series históricas de precios de activos financieros y construir el conjunto de datos necesario para el desarrollo experimental del trabajo.

Python es el lenguaje de programación elegido para el análisis de los datos y la implementación de la metodología propuesta. La manipulación y transformación de los datos se realizará mediante las librerías \textit{pandas} y \textit{NumPy}. Para el desarrollo de los modelos se emplearán librerías de aprendizaje profundo como \textit{PyTorch} y/o \textit{TensorFlow}. Asimismo, para la visualización de resultados se utilizará \textit{matplotlib}.

Para el control de versiones se ha empleado Git, utilizando la plataforma GitHub para el almacenamiento remoto del código, la documentación asociada y el seguimiento de los cambios realizados durante el desarrollo del trabajo.
